\section{Methods}
\label{sec:methods}

\begin{figure}[h]
  \centering
  \figmaybe[width=0.98\textwidth]{main/fig0_concept_overview.png}
  \caption{Overview of the jaxon differentiable pipeline.
    \textbf{(a) Forward-and-optimize loop.}  A per-contact
    stimulus $I_k(t)$ drives the superposed extracellular field
    $V_e = \sum_k I_k\,V_e^{(k)}$; the coupled-cable GPU solver
    integrates the membrane dynamics of every fiber in parallel and
    aggregates them into a fascicle-selectivity objective
    $\mathcal{L}$.  Because the whole pipeline is built on
    \texttt{jaxley}/JAX, the gradient
    $\partial\mathcal{L}/\partial\theta$ with respect to the stimulus
    is available end-to-end by reverse-mode autodiff (dashed loop).
    The solver is validated against \texttt{pyfibers}/NEURON on
    thresholds and conduction velocity.
    \textbf{(b) Anatomy to model.}  A histological vagus-nerve
    cross-section is segmented and embedded in the 12-contact
    MultiContact cuff finite-element model, from which the per-contact
    extracellular lead fields $V_e^{(k)}$ are computed.
    \textbf{(c) Fiber models and validation.}  The four
    implemented fiber models (myelinated MRG and Sweeney; unmyelinated
    Sundt and Rattay) and their strength--duration agreement with
    \texttt{pyfibers}/NEURON.
    \textbf{(d) Selectivity optimization on the cohort.}
    Two example nerves with the optimized cuff pattern (cathode blue,
    anode orange, inactive white, over the cuff's four angular columns
    $\times$ three axial rows) and fibers colored by recruitment
    (target fired / target silent / off-target fired / off-target
    silent); the peripheral target-fascicle cluster is tinted green.}
  \label{fig:overview}
\end{figure}

\subsection{Axon models}
\label{sec:methods-models}

jaxon implements four canonical peripheral nerve fiber models: the
myelinated McIntyre--Richardson--Grill (MRG) double-cable and Sweeney
single-cable axons, and the unmyelinated Sundt and Rattay
C-fibers~\cite{mcintyre2002,sweeney1987,sundt2015,rattay1989analysis}.
Each carries its published channel set and diameter-dependent geometry
(resting potentials $-80$\,mV for MRG/Sweeney, $-60$\,mV Sundt,
$-70$\,mV Rattay); the MRG model is a coupled $(V_i, V_{\mathrm{px}})$
double cable, the others single cables.  The governing equations and
all constants are given in the Supplementary Material; all four share
a common compartment-level interface (\texttt{FiberStatics},
\texttt{batch\_integrate\_m\_max}), so the library extends to further
models without solver changes.  Figure~\ref{fig:overview} summarizes
the forward-and-optimize pipeline, FE construction, validation and
cohort optimization.

\paragraph{Coupled solver and batched forward pass.}
The MRG double cable is integrated by a backward-Euler scheme in which
the capacitance-free periaxonal equation is eliminated by a Schur
complement and the resulting tridiagonal system solved by a
JAX-compatible Thomas sweep (Supplementary Material).  The
forward pass for one fiber is wrapped in \texttt{jax.lax.scan} over
time steps and returns, for optimization, only the per-compartment
maximum of the Na$^+$ activation gate $\mathbf{m}_{\mathrm{Na}}$ --- a
differentiable recruitment proxy --- at $\ge\!2$ orders of magnitude
lower memory than the full trace.  For $N$ fibers under the same
electrode, the per-fiber static parameters are stacked and a single
\texttt{jax.vmap} call compiles to one $N$-way XLA kernel, eliminating
the per-fiber Python overhead that dominates NEURON pipelines and
driving the throughput gain (Section~\ref{sec:results-scaling}).  We
integrate at $\Delta t = 5\,\mu\mathrm{s}$.

\subsection{Selectivity optimization}
\label{sec:methods-opt}

\paragraph{Loss function.}
For each fiber $f$ we form a differentiable activation proxy
$a_f \in [0,1]$ from the per-compartment maximum over time of the
Na$^+$ activation gate, $\mu_c = \max_t \mathbf{m}_{\mathrm{Na}}(t)|_c$.
At each of the two terminal nodes $\mu$ is passed through a sigmoid,
$\sigma\!\bigl((\mu - \tfrac12)/\tau\bigr)$ with temperature
$\tau = 0.1$, and $a_f$ is the smaller of the two ends --- so
$a_f \to 1$ when an action potential propagates to both ends and
$a_f \to 0$ otherwise.  The sigmoid smoothing keeps the gradient
informative away from threshold; its binary limit
($\tau\!\to\!0$) is $a_f = \min(\mu_{\mathrm{first}},
\mu_{\mathrm{last}})$, whose gradient vanishes except in a narrow
window at threshold.  Recruitment counts and the reported selectivity
index threshold $a_f$ at $0.5$.  We minimize the weighted-quality (WQ)
loss
\begin{equation}
  L_{\mathrm{WQ}} = \sum_f w_f \left[ y_f (1 - a_f) + (1 - y_f) a_f \right]
                 + \lambda \cdot \frac{1}{K} \sum_k I_k^2,
  \label{eq:wq-loss}
\end{equation}
where $y_f \in \{0,1\}$ is the target/off-target indicator
(constructed per nerve by the peripheral-cluster selector of
Section~\ref{sec:methods-duke}) and the $w_f$ are class-balancing
weights (each of the target and off-target populations carries total
weight $\tfrac12$, so neither dominates the imbalanced split).  The
second term is a small amplitude-energy regularizer that discourages
the trivial ``fire-everything'' optimum; we set $\lambda = 10^{-3}$ for
all results reported here, small enough to leave the single-diameter
selectivity outcomes unchanged.

\paragraph{Amplitude optimization.}
The design variable is the per-contact amplitude vector
$\mathbf{I} \in \mathbb{R}^K$ ($K=12$); each contact scales a fixed,
charge-balanced asymmetric-biphasic rectangular pulse ($0.5\,$ms
cathodic phase followed by a $2.0\,$ms anodic recharge at one-quarter
amplitude, so net charge is zero --- the clinical convention for
chronic implants).  All selectivity results use
Adam~\cite{kingma2015adam} with finite-difference gradients: each step
needs only $K{+}1$ forward passes and avoids back-propagating through
the long integration, so it runs on the full ${\sim}1000$-fiber
population --- the regime prior NEURON pipelines could not afford.  The
amplitude clip, learning rate and finite-difference step are listed in
Supplementary Table~\ref{tab:opt-hparams}; the loop runs for up to $90$
iterations and early-stops when the hard SI reaches $0.95$ (a
convergence floor: an SI within $0.05$ of the perfect-selectivity
ceiling, the resolution below which additional iterations no longer
change which fibers fire) or when the loss and hard SI both plateau.  On the histology cohort the trajectory is
warm-started by a discrete probe of physically meaningful spatial
patterns --- one bipolar and four single-column tripolar (cathode at
each column's $z$-middle contact with charge-balancing anodes) --- each
scored at eight magnitudes in both polarities ($80$ forward passes per
seed); the best initializes the trajectory, with the probe's zero
contacts frozen to preserve sparsity.  (jaxon also exposes exact
autodiff gradients, used for waveform-shape optimization; future work.)

\subsection{Subject-specific Duke FEM cohort}
\label{sec:methods-duke}

For the cohort-scale selectivity study (Section~\ref{sec:results-duke})
we use subject-specific finite-element
(FE) nerve models built from quantified vagus-nerve morphology of
swine~\cite{pelot2020swineSPARC} and human~\cite{pelot2021humanSPARC},
both released through the SPARC Portal.  The raw datasets provide,
per nerve, calibrated histological cross-sections with manually
segmented epineurium and per-fascicle endoneurium boundaries; the
human dataset additionally includes anti-claudin-1 staining that
resolves the perineurium thickness directly~\cite{pelot2023humanvagus}.

An automated pipeline (Golgi~\cite{haberbusch2026golgi}) converts each
segmented cross-section into a watertight 3-D multi-domain nerve-and-cuff
FE model --- nerve embedded in a concentric 12-contact cuff (3 axial
rows $\times$ 4 angular columns), saline film and muscle bath --- and
solves the linear electro-quasistatic problem
$\nabla \cdot (\sigma \nabla V) = 0$ in FEniCSx with literature
conductivities and a zero-thickness perineurium contact impedance,
giving a per-contact extracellular lead field $V_e^{(k)}$
(full construction, conductivities and the optimizer defaults in the
Supplementary Material).  The segmented anatomy, 3-D model and computed
fields for every cohort nerve are shown in Supplementary
Figs~\ref{fig:anat-human-sub-47-sam-2}--\ref{fig:anat-sub-9-sam-3}.

\paragraph{Fiber sampling, target selection and cohort gating.}
$1\,000$ straight axial fibers per nerve are seeded by area-weighted
rejection sampling inside the fascicles and the lead fields
$V_e^{(k)}$ sampled at every MRG compartment, giving a per-fiber tensor
${V_e^{(k)}}_{f,n}$ that the optimizer superposes linearly with the
amplitudes.  The target/off-target labels $y_f$ in
Eq.~(\ref{eq:wq-loss}) are set per nerve by a peripheral-cluster
selector that matches what the cuff can steer: a $90^\circ$ sector (the
cuff's angular resolution) is slid around the perimeter, and four
maximally-separated sectors carrying $\ge\!50$ peripheral fibers each
become the four target ``seeds'' optimized per nerve, with the
enclosed fascicles labeled target.  Of $36$ SPARC bundles processed,
$29$ nerves (18 swine, 11 human; 116 seeds) entered the study; the rest
were rejected for having $<\!3$ fascicles or insufficient FE-mesh
$z$-resolution between contact rows (full criteria in the Supplementary
Material).

\subsection{Validation protocol}
\label{sec:methods-validation}

For each model we drove jaxon and \texttt{pyfibers}-wrapped NEURON with
identical stimuli, at a matched fixed time step ($\Delta t = 5\,\mu$s;
NEURON integrated with its default backward-Euler method), and compared
three quantities against NEURON: the
strength--duration threshold (terminal node, $\mathrm{PW} \in
[0.02,1.0]\,\mathrm{ms}$), the conduction velocity over the published
diameter range, and the full center-node intracellular trace; agreement
to $<\!1\,\%$ on all three was required for acceptance, following the
established practice of validating open-source nerve-stimulation models
against NEURON~\cite{bucksot2021validation}.  All results
were produced on one NVIDIA A100 (40\,GB) with \texttt{jax} 0.6.2 /
\texttt{jaxley} 0.13.0 (CUDA 12.4) and pyfibers 0.8.5 (NEURON 8.2); the
full environment and reproduction code are available in the project
repository.

\subsection{Statistical analysis}
\label{sec:methods-stats}

Selectivity comparisons on the histology cohort are paired by nerve
$\times$ target-cluster seed, so we test them with the two-sided
Wilcoxon signed-rank test, which makes no Gaussian assumption on the
heavy-tailed, ceiling-bounded SI distributions.  For the
sparse-sampling analysis (Section~\ref{sec:results-sparse}) the
independent unit is the nerve: per-seed values are averaged within each
nerve before testing, so that nerves rather than nested seeds are the
units of replication.  The deployment penalty (dense-optimized SI minus
sparse-optimized transfer SI) is tested per species and sampling
strategy with two-sided Wilcoxon signed-rank tests, Holm-corrected
across the eight species $\times$ strategy cells; the species
difference in penalty is tested with the Mann--Whitney $U$ test, and its
anatomical driver --- the number of fibers per target fascicle --- with
Spearman rank correlation.  Seeds whose target cluster is degenerate
(target fraction $> 0.6$ --- spanning most of the nerve, so no selective
solution exists by construction; $5$ of $116$, all human) are excluded
from all analyses; low-SI outcomes on selectable targets are retained as
genuine geometric-ceiling cases rather than solver failures.  We
report medians with inter-quartile ranges alongside means throughout,
because the SI distributions are bounded at $1.0$ and right-skewed by
the geometric-ceiling subset, so the median is the more representative
central estimate.  Analyses use \texttt{scipy.stats}.
